% GENERATED FILE. Edit canonical lessons, then run npm run generate:books.
% canonical-source-hash: fnv1a64:82dfaaf0bb11b556
% canonical-lessons: ZH-C69-zuofan, ZH-W69-zuo-make, ZH-C69-zaofan, ZH-C69-wufan, ZH-C69-wanfan, ZH-C69-haochi, ZH-R69-review-64-69

\chapter{Meals}
\label{ch:zh-meals}
\hlchaptermodality{69}

\begin{chapteropening}
\textbf{By the end of this chapter:} \emph{I can name the meals of the day, say who cooks and say it tastes good.}

Complete the last lesson of chapter 69: I can name the meals of the day, say who cooks and say it tastes good.
\end{chapteropening}

\section[\texorpdfstring{\zh{做饭} (zuòfàn)}{做饭 (zuòfàn)}]{\zh{做饭} (zuòfàn) — to cook}

\label{lesson:ZH-C69-zuofan}



\begin{quote}
Before the new one: say the Mandarin for flower, then the Mandarin for the sun.
\end{quote}

\subsection*{You'll want to know: \zh{做饭}}
\textbf{\zh{做饭}} — \emph{zuòfàn} — \textquotedblleft{}to cook\textquotedblright{}.

\begin{quote}
\textbf{\zh{我妈妈喜欢做饭}} \emph{wǒ māma xǐhuan zuòfàn} — My mum loves cooking.
\end{quote}

The new character, in pieces: \zh{做} — \zh{亻} on the left, \zh{十} over \zh{口} in the middle, and on the right a four-stroke piece like a hand striking. The next lesson puts it in your hand.

\begin{cousinweb}
\textbf{\zh{做饭}} (\emph{zuòfàn}) is built from \zh{做} 'to do, to make' and \zh{饭} 'meal'. \zh{饭} is inherited from Middle Chinese \emph{bjonH} (Guangyun, in Baxter's transcription); \zh{做} has no reading in any dictionary we use, so its history is not settled.
\end{cousinweb}

\begin{grammarlens}[title={What you've built}]
One more word for everyday answers.
\end{grammarlens}

\subsection*{Your turn}
\begin{itemize}
\raggedright
  \item \emph{Say it:} \emph{zuòfàn}
  \item \emph{Say it:} \emph{zuòfàn}, once more
  \item \emph{Say it:} \emph{wǒ māma xǐhuan zuòfàn} — My mum loves cooking.
  \item \emph{Recall:} say the Mandarin for flower, then the Mandarin for the sun, then say \emph{zuòfàn} again
\end{itemize}

\begin{tcolorbox}[breakable,colback=teal!4,colframe=teal!35!black,title={Before you move on}]
What does \zh{做饭} mean? (\textquotedblleft{}To cook\textquotedblright{}.) Say it once more.
\end{tcolorbox}
\section[\texorpdfstring{\zh{做} (zuò)}{做 (zuò)}]{\zh{做} — the person, gǔ, and a striking hand}

\label{lesson:ZH-W69-zuo-make}



\begin{quote}
The word first, then the character inside it.

\begin{itemize}
\raggedright
  \item \emph{Recall:} say \emph{zuòfàn} — \textbf{R1}, one lesson back, before it decays
\end{itemize}

\begin{quote}
\zh{他} \zh{十} \zh{口}
\end{quote}

\emph{Read it:} each one aloud, with nothing said first — \emph{tā}, \emph{shí}, \emph{kǒu}
\end{quote}

\subsection*{Script — the person, gǔ, and a striking hand}
\hlblockfigure{\detokenize{figures/ZH-W69-zuo-make-filmstrip.pdf}}{How \zh{做} is written, stroke by stroke}

\begin{quote}
\zh{做}
\end{quote}

\textbf{Three pieces, left to right.} Count three, not eleven.

\textbf{\zh{亻}, two strokes}, as in \zh{你} and \zh{他}.

\textbf{The middle piece, five strokes, read gǔ}: \zh{十}, a horizontal and a vertical, standing on a small \zh{口}.

\textbf{The right piece, four strokes}: a short left-falling stroke, a short horizontal, a long left-falling stroke, and a long stroke falling to the right.

2 + 5 + 4 = eleven.

\subsection*{Your turn}
\begin{itemize}
\raggedright
  \item \emph{Write it:} \zh{做}, all eleven, with the model in front of you
  \item \emph{Cover:} write it again from memory, and count the pieces, not the strokes
\end{itemize}

\begin{tcolorbox}[breakable,colback=teal!4,colframe=teal!35!black,title={Before you move on}]
How many pieces in \zh{做}? (\textbf{Three}.) What sits in the middle? (\textbf{\zh{十} on \zh{口}}.) How does it end? (\textbf{With a stroke falling to the right}.)

Source: \href{https://raw.githubusercontent.com/chanind/hanzi-writer-data/68d10a4b21150cae5e1ebbd223eed289cf32d90c/data/\%E5\%81\%9A.json}{Hanzi Writer Data \zh{做}, pinned stroke paths and medians}.
\end{tcolorbox}
\section[\texorpdfstring{\zh{早饭} (zǎofàn)}{早饭 (zǎofàn)}]{\zh{早饭} (zǎofàn) — breakfast}

\label{lesson:ZH-C69-zaofan}



\begin{quote}
Before the new one: say the Mandarin for the sun, then the Mandarin for to cook.
\end{quote}

\subsection*{You'll want to know: \zh{早饭}}
\textbf{\zh{早饭}} — \emph{zǎofàn} — \textquotedblleft{}breakfast\textquotedblright{}.

\begin{quote}
\textbf{\zh{我七点吃早饭}} \emph{wǒ qī diǎn chī zǎofàn} — I have breakfast at seven.
\end{quote}

\begin{cousinweb}
\textbf{\zh{早饭}} (\emph{zǎofàn}) is built from zǎo 'morning' and \zh{饭} 'meal', both inherited from Middle Chinese, \emph{tsawX} and \emph{bjonH} (Guangyun, in Baxter's transcription).
\end{cousinweb}

\begin{grammarlens}[title={What you've built}]
One more word for everyday answers.
\end{grammarlens}

\subsection*{Your turn}
\begin{itemize}
\raggedright
  \item \emph{Say it:} \emph{zǎofàn}
  \item \emph{Say it:} \emph{zǎofàn}, once more
  \item \emph{Say it:} \emph{wǒ qī diǎn chī zǎofàn} — I have breakfast at seven.
  \item \emph{Recall:} say the Mandarin for the sun, then the Mandarin for to cook, then say \emph{zǎofàn} again
\end{itemize}

\begin{tcolorbox}[breakable,colback=teal!4,colframe=teal!35!black,title={Before you move on}]
What does \zh{早饭} mean? (\textquotedblleft{}Breakfast\textquotedblright{}.) Say it once more.
\end{tcolorbox}
\section[\texorpdfstring{\zh{午饭} (wǔfàn)}{午饭 (wǔfàn)}]{\zh{午饭} (wǔfàn) — lunch}

\label{lesson:ZH-C69-wufan}



\begin{quote}
Before the new one: say the Mandarin for to cook, then the Mandarin for breakfast.
\end{quote}

\subsection*{You'll want to know: \zh{午饭}}
\textbf{\zh{午饭}} — \emph{wǔfàn} — \textquotedblleft{}lunch\textquotedblright{}.

\begin{quote}
\textbf{\zh{你吃午饭了吗}} \emph{nǐ chī wǔfàn le ma} — Have you had lunch?
\end{quote}

\begin{cousinweb}
\textbf{\zh{午饭}} (\emph{wǔfàn}) is built from \zh{午} 'noon' and \zh{饭} 'meal', both inherited from Middle Chinese, \emph{nguX} and \emph{bjonH} (Guangyun, in Baxter's transcription).
\end{cousinweb}

\begin{grammarlens}[title={What you've built}]
One more word for everyday answers.
\end{grammarlens}

\subsection*{Your turn}
\begin{itemize}
\raggedright
  \item \emph{Say it:} \emph{wǔfàn}
  \item \emph{Say it:} \emph{wǔfàn}, once more
  \item \emph{Say it:} \emph{nǐ chī wǔfàn le ma} — Have you had lunch?
  \item \emph{Recall:} say the Mandarin for to cook, then the Mandarin for breakfast, then say \emph{wǔfàn} again
\end{itemize}

\begin{tcolorbox}[breakable,colback=teal!4,colframe=teal!35!black,title={Before you move on}]
What does \zh{午饭} mean? (\textquotedblleft{}Lunch\textquotedblright{}.) Say it once more.
\end{tcolorbox}
\section[\texorpdfstring{\zh{晚饭} (wǎnfàn)}{晚饭 (wǎnfàn)}]{\zh{晚饭} (wǎnfàn) — dinner; supper}

\label{lesson:ZH-C69-wanfan}



\begin{quote}
Before the new one: say the Mandarin for breakfast, then the Mandarin for lunch.
\end{quote}

\subsection*{You'll want to know: \zh{晚饭}}
\textbf{\zh{晚饭}} — \emph{wǎnfàn} — \textquotedblleft{}dinner; supper\textquotedblright{}.

\begin{quote}
\textbf{\zh{今天晚饭吃什么}} \emph{jīntiān wǎnfàn chī shénme} — What's for dinner tonight?
\end{quote}

\begin{cousinweb}
\textbf{\zh{晚饭}} (\emph{wǎnfàn}) is built from wǎn 'evening' and \zh{饭} 'meal', both inherited from Middle Chinese, \emph{mjonX} and \emph{bjonH} (Guangyun, in Baxter's transcription).
\end{cousinweb}

\begin{grammarlens}[title={What you've built}]
One more word for everyday answers.
\end{grammarlens}

\subsection*{Your turn}
\begin{itemize}
\raggedright
  \item \emph{Say it:} \emph{wǎnfàn}
  \item \emph{Say it:} \emph{wǎnfàn}, once more
  \item \emph{Say it:} \emph{jīntiān wǎnfàn chī shénme} — What's for dinner tonight?
  \item \emph{Recall:} say the Mandarin for breakfast, then the Mandarin for lunch, then say \emph{wǎnfàn} again
\end{itemize}

\begin{tcolorbox}[breakable,colback=teal!4,colframe=teal!35!black,title={Before you move on}]
What does \zh{晚饭} mean? (\textquotedblleft{}Dinner; supper\textquotedblright{}.) Say it once more.
\end{tcolorbox}
\section[\texorpdfstring{\zh{好吃} (hǎochī)}{好吃 (hǎochī)}]{\zh{好吃} (hǎochī) — tasty; good to eat}

\label{lesson:ZH-C69-haochi}



\begin{quote}
Before the new one: say the Mandarin for lunch, then the Mandarin for dinner.
\end{quote}

\subsection*{You'll want to know: \zh{好吃}}
\textbf{\zh{好吃}} — \emph{hǎochī} — \textquotedblleft{}tasty; good to eat\textquotedblright{}.

\begin{quote}
\textbf{\zh{我妈妈做饭很好吃}} \emph{wǒ māma zuòfàn hěn hǎochī} — My mum's cooking is delicious.
\end{quote}

\begin{cousinweb}
\textbf{\zh{好吃}} (\emph{hǎochī}) is built from \zh{好} 'good' and \zh{吃} 'to eat': good to eat. \zh{好} is inherited from Middle Chinese \emph{xawX} (Guangyun, in Baxter's transcription); \zh{吃} is probably the \emph{khek} the dictionary writes with another character and glosses 'to eat', though chī is not the sound \emph{khek} would regularly give.
\end{cousinweb}

\begin{grammarlens}[title={What you've built}]
One more word for everyday answers.
\end{grammarlens}

\subsection*{Your turn}
\begin{itemize}
\raggedright
  \item \emph{Say it:} \emph{hǎochī}
  \item \emph{Say it:} \emph{hǎochī}, once more
  \item \emph{Say it:} \emph{wǒ māma zuòfàn hěn hǎochī} — My mum's cooking is delicious.
  \item \emph{Recall:} say the Mandarin for lunch, then the Mandarin for dinner, then say \emph{hǎochī} again
\end{itemize}

\begin{tcolorbox}[breakable,colback=teal!4,colframe=teal!35!black,title={Before you move on}]
What does \zh{好吃} mean? (\textquotedblleft{}Tasty; good to eat\textquotedblright{}.) Say it once more.
\end{tcolorbox}
\section[(dialogue)]{Looking back: from Asking for Help to Meals}

\label{lesson:ZH-R69-review-64-69}



\begin{quote}
Say the Mandarin for sorry and to help.
\end{quote}

\subsection*{Your turn}
\begin{itemize}
\raggedright
  \item \emph{Say it:} sorry, to help, to do a favour, to borrow, to give back
  \item \emph{Say it:} door, may I ask, excuse me, to open, to shut
  \item \emph{Say it:} weather, rain, to rain, snow, wind
  \item \emph{Say it:} horse, dog, cat, bird, fish
  \item \emph{Say it:} mountain, star, grass, flower, the sun
  \item \emph{Say it:} to cook, breakfast, lunch, dinner, tasty
  \item \emph{Read it:} \zh{起} \zh{门} \zh{问} \zh{气} \zh{雨} \zh{马} \zh{山} \zh{做} — each one aloud, then a word that holds it
\end{itemize}

\begin{tcolorbox}[breakable,colback=teal!4,colframe=teal!35!black,title={Before you move on}]
Sorry? (\textbf{\zh{对不起}}, \emph{duìbuqǐ}.) To help? (\emph{bāng}.) To do a favour? (\emph{bāngmáng}.) To borrow? (\emph{jiè}.) To give back? (\emph{huán}.) Door? (\textbf{\zh{门}}, \emph{mén}.) May I ask? (\textbf{\zh{请问}}, \emph{qǐngwèn}.) Excuse me? (\emph{bù hǎoyìsi}.) To open? (\emph{kāi}.) To shut? (\emph{guān}.) Weather? (\textbf{\zh{天气}}, \emph{tiānqì}.) Rain? (\textbf{\zh{雨}}, \emph{yǔ}.) To rain? (\emph{xiàyǔ}.) Snow? (\emph{xuě}.) Wind? (\emph{fēng}.) Horse? (\textbf{\zh{马}}, \emph{mǎ}.) Dog? (\emph{gǒu}.) Cat? (\emph{māo}.) Bird? (\emph{niǎo}.) Fish? (\emph{yú}.) Mountain? (\textbf{\zh{山}}, \emph{shān}.) Star? (\textbf{\zh{星星}}, \emph{xīngxing}.) Grass? (\emph{cǎo}.) Flower? (\emph{huā}.) The sun? (\emph{tàiyang}.) To cook? (\textbf{\zh{做饭}}, \emph{zuòfàn}.) Breakfast? (\textbf{\zh{早饭}}, \emph{zǎofàn}.) Lunch? (\textbf{\zh{午饭}}, \emph{wǔfàn}.) Dinner? (\textbf{\zh{晚饭}}, \emph{wǎnfàn}.) Tasty? (\textbf{\zh{好吃}}, \emph{hǎochī}.)
\end{tcolorbox}
